% Source: https://artofproblemsolving.com/community/c3h1569195_painters_and_apartments
% Stable identity: aops-topic-1569195-q01
Three painters working $3$ days give $3\cdot3=9$ painter-days of work and paint $3$ apartments, so one apartment takes $3$ painter-days (one painter paints $\dfrac{1}{3}$ of an apartment per day). The number of apartments painted is
\[
\dfrac{\text{painters}\times\text{days}}{3}.
\]
\begin{enumerate}
\item[(a)] $\dfrac{6\cdot6}{3}=\boxed{12}$ apartments.
\item[(b)] $9$ apartments need $27$ painter-days, so $9$ painters need $\dfrac{27}{9}=\boxed{3}$ days.
\item[(c)] $12$ apartments need $36$ painter-days, so in $12$ days we need $\dfrac{36}{12}=\boxed{3}$ painters.
\end{enumerate}
